% Part: second-order-logic
% Chapter: metatheory
% Section: second-order-arithmetic

\documentclass[../../../include/open-logic-section]{subfiles}

\begin{document}

\olfileid{sol}{met}{spa}

\olsection{Aritmetika Orde Kapindho}

Elinga yen teori $\Th{PA}$ saka aritmetika Peano ngemot
wolung aksioma saka~$\Th{Q}$,
\begin{align*}
& \lforall[x][\eq/[x'][\Obj 0]]\\
& \lforall[x][\lforall[y][(\eq[x'][y'] \lif \eq[x][y])]]\\
& \lforall[x][(\eq[x][\Obj 0] \lor \lexists[y][\eq[x][y']])]\\ 
& \lforall[x][\eq[(x + \Obj 0)][x]]\\
& \lforall[x][\lforall[y][\eq[(x + y')][(x + y)']]]\\
& \lforall[x][\eq[(x \times \Obj 0)][\Obj 0]]\\
& \lforall[x][\lforall[y][\eq[(x \times y')][((x \times y) + x)]]]\\
& \lforall[x][\lforall[y][(x < y \liff \lexists[z][\eq[(z' + x)][y]])]]\\
\intertext{ditambah kabeh ukara kanthi wujud}
& (!A(\Obj 0) \land \lforall[x][(!A(x) \lif !A(x'))]) \lif \lforall[x][!A(x)].
\end{align*}
Kang pungkasan iku ``skema,'' yaiku pola kang ngasilake
!!{sentence}s basa aritmetika kang cacahe tanpa wates,
siji kanggo saben !!{formula}~$!A(x)$. Skema iki diarani
\emph{skema aksioma induksi} (orde kapisan).
Ing aritmetika Peano \emph{orde kapindho}~$\Th{PA^2}$,
induksi bisa dinyatakake nganggo siji ukara.
$\Th{PA^2}$ ngemot wolung aksioma kapisan ing ndhuwur
ditambah \emph{aksioma induksi} (orde kapindho):
\[
\lforall[X][((X(\Obj 0) \land \lforall[x][(X(x) \lif X(x'))]) \lif \lforall[x][X(x)])].
\]
Iki nyatakake yen subhimpunan~$X$ saka !!{domain}
ngemot~$\Assign{\Obj{0}}{M}$ lan, kanggo saben~$x \in \Domain{M}$
kang kalebu ing subhimpunan kasebut, uga
ngemot~$\Assign{\prime}{M}(x)$ (yaiku ``tutup tumrap
panerus''), mula subhimpunan iku ngemot sakabehe unsur
ing !!{domain} (yaiku $X = \Domain{M})$.

Aksioma induksi njamin yen saben !!{structure} kang nyukupi
aksioma iku mung ngemot !!{element}s saka~$\Domain{M}$
kang pancen diwajibake dening aksioma, yaiku nilai
saka~$\num{n}$ kanggo $n \in \Nat$. Model
$\Th{PA^2}$ ora ngemot wilangan nonstandar.

\begin{thm}
\ollabel{thm:sol-pa-standard}
Yen $\Sat{M}{\Th{PA^2}}$, mula $\Domain{M} =
\Setabs{\Value{\num{n}}{M}}{n \in \Nat}$.
\end{thm}

\begin{proof}
Upamana $N = \Setabs{\Value{\num{n}}{M}}{n \in \Nat}$, lan
upamana $\Sat{M}{\Th{PA^2}}$. Mesthi wae, kanggo saben
$n \in \Nat$, $\Value{\num{n}}{M} \in \Domain{M}$,
mula $N \subseteq \Domain{M}$.

Saiki buktekake inklusi arah kosok baline. Gatekna
pangaji variabel $s$ kanthi $s(X) = N$. Miturut anggepan,
\begin{align*}
\Sat{M}{\lforall[X][((X(\Obj 0) \land \lforall[x][(X(x)
      \lif X(x'))]) \lif \lforall[x][X(x)])]}, & \text{ mula}\\
\Sat{M}{(X(\Obj 0) \land \lforall[x][(X(x) \lif X(x'))]) \lif
  \lforall[x][X(x)]}[s]. & 
\end{align*}
Gatekna anteseden saka implikasi iki.
$\Value{\Obj{0}}{M} \in N$, mula
$\Sat{M}{X(\Obj{0})}[s]$. Konjungsi kapindho,
$\lforall[x][(X(x) \lif X(x'))]$, uga dicukupi.
Upamana $x \in N$. Miturut definisi~$N$,
$x = \Value{\num{n}}{M}$ kanggo sawijining~$n$.
Mula $\Assign{\prime}{M}(x) = \Value{\num{n+1}}{M} \in N$.
Dadi, $\Assign{\prime}{M}(x) \in N$.

Kita nduweni
$\Sat{M}{X(\Obj 0) \land \lforall[x][(X(x) \lif X(x'))]}[s]$.
Mula $\Sat{M}{\lforall[x][X(x)]}[s]$. Iki ateges
kanggo saben $x \in \Domain{M}$, kita nduweni
$x \in s(X) = N$. Dadi, $\Domain{M} \subseteq N$.
\end{proof}

\begin{cor}
\ollabel{cor:sol-pa-categorical}
Sembarang rong model saka~$\Th{PA^2}$ iku isomorfik.
\end{cor}

\begin{proof}
Miturut \olref{thm:sol-pa-standard}, domain saben model saka
$\Th{PA^2}$ kapenuhan dening $\Value{\num{n}}{M}$.
Saben model kaya mangkono uga model saka~$\Th{Q}$.
Miturut \olref[mod][mar][stm]{prop:thq-standard}, saben
model kasebut standar, yaiku isomorfik karo~$\Struct{N}$.
\end{proof}

Ing ndhuwur kita netepake $\Th{PA^2}$ minangka teori kang ngemot
wolung aksioma aritmetika kapisan ditambah aksioma induksi orde
kapindho. Sejatine, amarga daya ungkap logika orde kapindho,
aritmetika Peano orde kapindho mung mbutuhake \emph{rong aksioma
kapisan} saka aksioma aritmetika mau, ditambah induksi.

\begin{prop}\ollabel{prop:sol-pa-definable}
Upamana $\Th{PA^{2\dagger}}$ iku teori orde kapindho kang ngemot
rong aksioma aritmetika kapisan (aksioma panerus) lan aksioma
induksi orde kapindho. Mula $\le$, $+$, lan $\times$ bisa
didefinisikake ing $\Th{PA^{2\dagger}}$.
\end{prop}

\begin{proof}
Kanggo nuduhake yen $\le$ bisa didefinisikake, kita kudu
nemokake formula~$!A_\le(x, y)$ saengga
$\Sat{N}{!A_\le(\num n, \num m)}$ yen lan mung yen
$n \le m$. Gatekna formula
\[
!B(x, Y) \ident Y(x) \land \lforall[y][(Y(y) \lif Y(y'))]
\]
Cetha yen $!B(\num n, Y)$ dicukupi dening sawijining
himpunan $Y \subseteq \Nat$, mula
$\Setabs{m}{n \le m} \subseteq Y$. Kosok baline ora kudu
bener, amarga himpunan kang ngemot buntut iki durung mesthi
tutup tumrap panerus ing unsur liyane. Nanging buntut iku
dhewe ngemot titik wiwitane lan tutup tumrap panerus. Cathetan koreksi sumber OLPL-806: ngemot kabeh unsur wiwit n mung akibat saka klausa B, dudu syarat kang ekuivalen; himpunan buntut iku dadi saksi kanggo arah kosok baline ing definisi universal. Mula
kita bisa njupuk $!A_\le(x, y) \ident
\lforall[Y][(!B(x, Y) \lif Y(y))]$.

Kanggo nuduhake yen tambah bisa didefinisikake, gatekna
$k+l=m$ yen lan mung yen ana fungsi $u$ saengga
$u(0)=k$, $u(n')=u(n)'$ kanggo saben $n$, lan
$m = u(l)$. Kita bisa nggunakake ekuivalensi iki kanggo
netepake tambah ing $\Th{PA^{2\dagger}}$ nganggo formula
ing ngisor iki:
\[
!A_+(x,y,z) \ident \lexists[u][(u(\Obj 0)=x \land \lforall[w][u(x')=u
 (x)'] \land u(y)=z)]
\] 
Cathetan koreksi sumber OLPL-807: kuantor kanggo w ing rumus kang dicithak ora dienggo ing klausa rekurensi, kang malah migunakake x. Kanggo iterasi panerus, klausa kudu lumaku kanggo saben w lan mbandhingake nilai u ing panerus w karo panerus nilai u ing w. Sawise klausa mau dibenerake, cetha yen $\Sat{N}{!A_+(\num k, \num l, \num m)}$
yen lan mung yen $k+l=m$.
\end{proof}

\begin{prob}
Rampungna bukti \olref[sol][met][spa]{prop:sol-pa-definable}.
\end{prob}

\end{document}
